\documentclass{situationist-preprint}

\title{Virtue Is Not Sufficient for Happiness}
\msid{ETH-11-a}
\msfamily{ETH-11}
\msdiscipline{Ethics}
\mstype{Principal manuscript}
\mskeywords{virtue; happiness; eudaimonia; Stoicism; Aristotle; Priam; goods of fortune; preferred indifferents}

\begin{document}
\maketitle

\begin{abstract}
The Stoics held that virtue is sufficient for happiness: the virtuous person is happy on the rack, and the goods of fortune contribute nothing to how well her life goes. Aristotle held that virtue is necessary for happiness but that a virtuous person who meets with the misfortunes of Priam cannot be called happy. We defend Aristotle. Happiness is a complete life that goes well for the one who lives it, and such a life includes goods, among them relationships, achievement and contact with reality, that are good for their possessor and that fortune can destroy. Virtue is necessary because without the virtues these goods cannot be had in the form that happiness requires: friendship needs justice and good faith, achievement needs perseverance and practical wisdom. Virtue is not sufficient because the goods can be lost by a virtuous agent through no fault of her own, and Priam, virtuous and bereaved of his sons and his city, is such an agent. We argue that the Stoic response, on which such losses are of preferred indifferents and leave happiness untouched, is unstable: a thing one has reason to prefer for one's own sake is good for one.
\end{abstract}

\section{The problem}

Priam, king of Troy, lived a just and pious life, had fifty sons, and ruled a prosperous city. In old age he saw most of his sons killed, among them Hector, whose body he had to beg from the man who killed him; his city was sacked, and he was killed at the altar of Zeus. \citet[I.9, 1100a5--9]{AristotleNE} uses him as the example of a man of whom no one would say that he was happy, though he was good. A virtuous person who meets with such fortunes is never miserable, since he will never do what is hateful and mean; but he is not blessed either \citep[I.10, 1100b33--1101a8]{AristotleNE}. Happiness, for Aristotle, is activity of the soul in accordance with virtue, in a complete life, and adequately furnished with external goods \citep[I.7--8, I.10]{AristotleNE}.

The Stoics denied this. Virtue is the only good, vice the only evil, and everything else (health, wealth, family, country, life itself) is indifferent. Some indifferents are \enquote{preferred} and to be selected when nothing prevents it, but their presence or absence makes no difference to happiness, which consists in virtue alone \citep[III]{CiceroFin}; the wise man is happy even under torture \citep[V]{CiceroTusc}. The thesis is not an extravagance but the conclusion of an argument: if the goods of fortune could make the difference between happiness and its absence, happiness would be at the mercy of chance, and the good person would not be in control of whether her life goes well \citep{LongSedley1987,Annas1993}.

The modern discussion has largely reaffirmed Aristotle's side of the question, with \citet{Cooper1985} on the role of external goods as the resources for virtuous activity, \citet{Nussbaum1986} on the fragility of the goods that make up a good human life, and \citet{Russell2012} on the embodied conception of happiness; but the Stoic position retains defenders and its argument is seldom met directly. We state the case against it as an argument, using the account of well-being defended in ETH-05 and the plural account of the good of ETH-01.

\section{Framework}

\begin{definition}[Happiness]\label{d:happy}
A person is \emph{happy} if her life, taken as a whole, goes well to a high degree for her: its level of well-being is high.
\end{definition}

Definition~\ref{d:happy} uses \enquote{happiness} for \emph{eudaimonia}, not for a feeling. It is neutral between Aristotle and the Stoics, who agree that happiness is the highest good for a person, the end for the sake of which everything else is chosen, and disagree about what it consists in.

\begin{definition}[Virtue]\label{d:virtue}
A person is \emph{virtuous} if she has the virtues of character, such as justice, courage, temperance and good faith, together with practical wisdom, to the degree attainable by a good human being.
\end{definition}

The last clause sets aside the Stoic sage, whose virtue is perfect and who, on the Stoics' own account, has rarely if ever existed. The question is whether the virtue of a good human being is sufficient for her happiness.

\begin{principle}[Goods of a life]\label{p:goods}
The well-being of a life includes, besides pleasure, the goods of relationships, achievement and contact with reality: other things equal, a life with more of these goods goes better for the one who lives it.
\end{principle}

Principle~\ref{p:goods} combines two results. ETH-05 shows, by comparing lives alike in experience, that well-being includes contact with reality: a life whose friendships and achievements are real goes better for its subject than an experiential twin whose friendships and achievements are illusory. ETH-01 identifies relationships and achievement as irreducible components of the good. We take the components of the good that are realized in a single life to be goods for its subject.

\begin{principle}[Preference and good]\label{p:pref}
If a person has reason, for her own sake and for nothing beyond it, to prefer that $x$ rather than $y$ be part of her life, then $x$ is better for her than $y$.
\end{principle}

Principle~\ref{p:pref} links reasons for self-regarding preference to well-being. It is what makes the notion of what is good \emph{for} a person intelligible: the good for a person is what there is reason to want for her sake \citep{Darwall2002}.

\section{The argument}

\begin{lemma}[Necessity]\label{l:nec}
Virtue is necessary for happiness.
\end{lemma}

\begin{proof}
By Principle~\ref{p:goods}, well-being includes relationships, achievement and contact with reality, and a life high in well-being (Definition~\ref{d:happy}) has a high level of these goods taken together. Each of them, in the measure that happiness requires, depends on the virtues. Friendship in the full sense is a relationship of mutual goodwill between people who wish each other well for each other's sake, and it requires justice and good faith: the unjust and the faithless can have associations of use and pleasure, but these dissolve when the use or pleasure ends \citep[VIII.3--4]{AristotleNE}. Achievement of any significance requires perseverance in the face of difficulty, which is a form of courage, and the judgment to choose worthwhile ends and suitable means, which is practical wisdom. Contact with reality, in the sense of ETH-05, requires that one's relationships be what one takes them to be, which a person who deceives and manipulates others cannot secure, since she gives those around her reason to deceive her in turn. A life lacking the virtues therefore lacks these goods in the measure happiness requires, and is not happy.
\end{proof}

The argument does not show that every virtuous person has these goods, which is the point of what follows, nor that a vicious person has none of them. It shows that their possession in the measure happiness requires depends on dispositions of character, so that without the virtues they are at best partial and precarious.

\begin{assumption}[Priam]\label{a:priam}
Priam was virtuous, and in his last years he lost, through no fault of his own, most of his sons, his city and his life.
\end{assumption}

The assumption is about the Priam of the tradition that Aristotle invokes; whether a historical king of Troy had these properties is irrelevant. It suffices that such a life is possible, and it plainly is: virtuous people lose their children to war and disease and their communities to conquest.

\begin{theorem}\label{t:insuff}
Virtue is not sufficient for happiness.
\end{theorem}

\begin{proof}
By Assumption~\ref{a:priam}, Priam was virtuous. Compare his life with an alternative life $L$, alike in his virtue and in everything up to the fall of Troy, in which Troy stands and his sons survive him. In $L$, Priam's relationships with his sons continue, his achievement as a ruler endures, and his contact with reality is undiminished. By Principle~\ref{p:goods}, $L$ goes better for Priam than his actual life. Moreover, what he lost is great: he lost most of what made up the goods of his life, at its end, when nothing could replace them. Taking his life as a whole, his well-being was far below what it would have been in $L$, and not high: his life did not go well for him to a high degree. By Definition~\ref{d:happy}, he was not happy. Since he was virtuous, virtue is not sufficient for happiness.
\end{proof}

\begin{corollary}\label{c:main}
Virtue is necessary but not sufficient for happiness.
\end{corollary}

\begin{proof}
Lemma~\ref{l:nec} and Theorem~\ref{t:insuff}.
\end{proof}

Corollary~\ref{c:main} preserves Aristotle's further claim that the virtuous person is never wretched. Her virtue is itself part of how her life goes, and it secures her against the worst evils, those she would do herself; it does not secure her against those that are done to her.

\section{Objections and replies}

\begin{objection}[Indifferents]\label{o:indiff}
The Stoic denies Principle~\ref{p:goods}. Sons and cities are preferred indifferents: Priam had reason to select them while he could, but their loss is not the loss of any good, and his happiness, which consists in virtue alone, is undiminished.
\end{objection}

\begin{reply}
The Stoic grants that Priam had reason to prefer that his sons live, and not merely as means to his virtuous activity, since his virtue would have been exercised in grief as fully as in prosperity. His reason was a reason to prefer their lives, for his own sake among others: the life of a father whose sons survive him is the life he had reason to want. By Principle~\ref{p:pref}, that life is better for him. The Stoic must therefore hold that there is reason to prefer something for one's own sake that is not better for one. The position is stable only if \enquote{preferred} and \enquote{better for one} come apart, and they do so only if \enquote{good for one} is redefined as \enquote{virtuous}, at which point the sufficiency thesis is true by stipulation and no longer conflicts with Aristotle's.
\end{reply}

\begin{objection}[Control]\label{o:control}
If happiness depends on fortune, the good person is not in control of whether her life goes well, and happiness becomes a matter of luck. This is unacceptable for a conception of the highest good.
\end{objection}

\begin{reply}
The premise that the highest good must be under the agent's control is what is in question, and it is not supported by reflection on what we want for those we love. We do not wish for our children only that they be good; we wish that their goodness not be met by the loss of everything they love. That happiness is subject to luck is a fact about human life, which consists in part in relationships with others and in projects in a world that does not always cooperate. What is in the agent's control is her virtue, and Lemma~\ref{l:nec} gives it the place it deserves: without it there is no happiness at all.
\end{reply}

\begin{objection}[A whole life]\label{o:whole}
Priam was happy for most of his life. If happiness belongs to a life as a whole, a few bad years at the end cannot undo decades of it.
\end{objection}

\begin{reply}
The objection treats a life's well-being as a sum of momentary levels. But the shape of a life matters to how well it goes: a life that ends in the destruction of everything its subject built and loved goes worse than one with the same momentary goods in reverse order \citep{Velleman1991}. Priam's loss falls on the goods of the whole life, his sons and his city, not only on its final years. This is Solon's point, which Aristotle discusses at length: one should call no one happy until one sees the end \citep[I.10]{AristotleNE}.
\end{reply}

\section{Conclusion}

Virtue is necessary for happiness, because the goods that make a life go well for its subject (real relationships, achievement, contact with reality) cannot be had in the measure happiness requires without the virtues. It is not sufficient, because those goods can be destroyed by fortune in the life of a virtuous person. Priam was virtuous and not happy, and the Stoic attempt to save the sufficiency thesis by calling his losses indifferent requires that one can have reason to prefer for one's own sake what is not good for one.

\begin{thebibliography}{99}\small
\bibitem[Annas(1993)]{Annas1993} Annas, J. (1993). \emph{The Morality of Happiness}. Oxford University Press.
\bibitem[Aristotle(NE)]{AristotleNE} Aristotle. \emph{Nicomachean Ethics}. Cited by book, chapter and Bekker page.
\bibitem[Cicero(Fin.)]{CiceroFin} Cicero. \emph{De Finibus Bonorum et Malorum}. Cited by book.
\bibitem[Cicero(Tusc.)]{CiceroTusc} Cicero. \emph{Tusculan Disputations}. Cited by book.
\bibitem[Cooper(1985)]{Cooper1985} Cooper, J.~M. (1985). Aristotle on the goods of fortune. \emph{The Philosophical Review}, 94(2), 173--196.
\bibitem[Darwall(2002)]{Darwall2002} Darwall, S. (2002). \emph{Welfare and Rational Care}. Princeton University Press.
\bibitem[Long and Sedley(1987)]{LongSedley1987} Long, A.~A. and Sedley, D.~N. (1987). \emph{The Hellenistic Philosophers}. Cambridge University Press.
\bibitem[Nussbaum(1986)]{Nussbaum1986} Nussbaum, M.~C. (1986). \emph{The Fragility of Goodness: Luck and Ethics in Greek Tragedy and Philosophy}. Cambridge University Press.
\bibitem[Russell(2012)]{Russell2012} Russell, D.~C. (2012). \emph{Happiness for Humans}. Oxford University Press.
\bibitem[Velleman(1991)]{Velleman1991} Velleman, J.~D. (1991). Well-being and time. \emph{Pacific Philosophical Quarterly}, 72(1), 48--77.
\end{thebibliography}

\end{document}
